\documentclass[11pt,letterpaper]{article}
\usepackage{technical-design-spec}
\setDocTitle{CI/CD Pipeline Architecture Review Standard}
\setDocSubtitle{Design review and release-readiness quality gates}
\setDocVersion{1.0}
\setDocStatus{Technical education / proposed reference}
\setDocOwner{CI/CD Engineering}
\setDocDate{30 October 2026}
\setDocSummary{Supplementary engineering practice}
\begin{document}
\makedoctitle
\makedocsummary
\tableofcontents
\clearpage
\section{Review scope}
Assess proposed CI architecture, not only a successful screenshot. Scoring weight is illustrative; reviewers may adapt it through a recorded change.

\section{Review criteria}
\begin{itemize}
\item Source-to-artifact identity and dependency lock integrity.
\item PR event and workflow injection defense.
\item Permissions, OIDC and deployment account isolation.
\item Deterministic local/CI parity and correct test staging.
\item Database change compatibility and recovery.
\item SAST/secret/dependency gate effectiveness and exception lifecycle.
\item Artifact digest verification, attestation and promotion.
\item Metrics quality, on-call ownership and feedback routing.
\end{itemize}
\section{Disposition}
Each criterion receives evidence, severity, reviewer, finding ID, target date and pass/conditional/fail result. A failing critical trust-boundary criterion blocks approval. Conditions must be limited in scope and time and approved by accountable owners.

\section{Artifacts}
\begin{itemize}
\item Diagram and decision register
\item Negative access test results
\item Release dossier and rollback game-day results
\item Completed printable scorecard and tracked findings
\end{itemize}
\end{document}
